\begin{frame}
  \frametitle{Ackroyd-Riyait Problems}
  % describes the Ackroyd-Riyait problem
  The problems of interest are the \glspl{arp}, which are a set of non-multiplying, time-independent, one-speed 2D neutron transport models that were created to demonstrate a numerical artifact called \textit{ray effects}~\cite{ACKROYD19891}.
  These appear when solving low-scattering (non-diffusive) problems using discrete ordinates methods.
  In this project, we are primarily concerned with the \gls{arp1}.

  \begin{figure}
  \centering
  \begin{tikzpicture}[scale=0.25]
    \draw[fill=MaterialColor,very thick,draw=MaterialColor] (0,0) rectangle (10,10);
    \draw[fill=VoidColor,very thick,draw=VoidColor] (0,0) rectangle (5,5);
    \draw[fill=SourceColor,very thick,draw=SourceColor] (0,0) rectangle (1.25,1.25);
    \draw[fill=none, very thick] (0,0) rectangle (10,10);

    \node[draw=none, fill=none, rotate=90] at (-1,5) {Reflective Boundary};
    \node[draw=none, fill=none, rotate=0] at (5,-1) {Reflective Boundary};
    \node[draw=none, fill=none, rotate=-90] at (11,5) {Void Boundary};
    \node[draw=none, fill=none, rotate=0] at (5,11) {Void Boundary};

    \draw[fill=SourceColor,very thick,align=left] (15,7.5) rectangle ++(1.5,1.5);
    \node[anchor=west] at (16.7,8.25) {= Source};

    \draw[fill=VoidColor,very thick,align=left] (15,5.0) rectangle ++(1.5,1.5);
    \node[anchor=west] at (16.7,5.75) {= Void};

    \draw[fill=MaterialColor,very thick,align=left] (15,2.5) rectangle ++(1.5,1.5);
    \node[anchor=west] at (16.7,3.25) {= Absorbing/Scattering Material};
  \end{tikzpicture}
  \caption{Schematic of \gls{arp1}.}
\end{figure}
\end{frame}

\begin{frame}
  \frametitle{Ackroyd-Riyait Problem 1 Material Properties}
  \gls{arp1} is defined by constant-value properties seen in \cref{tab:arp1-props}.
  \begin{table}
    \begin{tabular}{llll}
      \toprule
      Region         & Source (n/cm\textsuperscript{2}-s) & $\Sigma_t$ (1/cm) & $\Sigma_s$ (1/cm) \\
      \midrule
      Source         & 6.4                                & 0.8 (0.2)         & 0.0 (0.19)        \\
      Void           & 0                                  & 0                 & 0                 \\
      Abs/Scattering & 0                                  & 0.8 (0.2)         & 0.0 (0.19)        \\
      \bottomrule
    \end{tabular}
    \caption{Constant-value material properties for \gls{arp1}. Values in parentheses indicate values used for a scattering problem.}
    \label{tab:arp1-props}
  \end{table}
\end{frame}
